\begin{defi}[Exponential-kernel Hawkes process]\label{def.exponentialHawkes}
 A $\numEventTypes$-dimensional Hawkes process $\multiCountingProc$ is called an
 \emph{exponential-kernel Hawkes process} if its kernels are
 \begin{equation}\label{eq.exponentialKernel}
  \hawkesKernel\subscriptee(t) = \excitation\subscriptee \, e^{-\decay t},
  \qquad t \geq 0,
 \end{equation}
 with $\excitation\subscriptee \geq 0$ and a single decay rate $\decay > 0$
 shared by every pair of types.
 We call $\excitation$ the excitation matrix and $\decay$ the decay.
 For each type $e$ we call
 \begin{equation}\label{eq.decayedCounts}
  \decayedCounts_e(t) := \sum_{j \, : \, \arrivalTimes_{j} < t}
  e^{-\decay (t - \arrivalTimes_{j})}
 \end{equation}
 the decayed counts.
\end{defi}

From \eqref{eq.branchingMatrix},
\begin{equation}\label{eq.GammaIsAOverBeta}
 \branchingMatrix = \frac{\excitation}{\decay},
\end{equation}
so the branching matrix is dimensionless while the excitation matrix is a rate.

\begin{remark}\label{remark.countsAgainstRates}
 Notice that $\sum_e \branchingMatrix\subscriptee$ and $\sum_e \excitation\subscriptee$
 are different quantities, differing by $\decay$:
 the first is the expected number of direct offspring of one event of type $\eone$,
 the second is the jump in the total intensity that the same event produces,
 measured per unit of time.
\end{remark}


The intensity of an exponential-kernel Hawkes process can be written as
\begin{equation}\label{eq.intensityFromState}
 \intensity[ ](t) = \baseIntensity + \excitation \decayedCounts(t).
\end{equation}
Notice that it is the strict inequality in the sum of \eqref{eq.decayedCounts} that makes
$\decayedCounts$ left-continuous, and therefore makes \eqref{eq.intensityFromState}
predictable, as Definition \ref{def.compensator} requires of an intensity.

Between events $\decayedCounts$ solves $d\decayedCounts = -\decay \decayedCounts \, dt$,
and at an event of type $e$ it gains the unit vector $\unitVector_e$, so that
\begin{equation}\label{eq.stateDynamics}
 d\decayedCounts(t) = -\decay \decayedCounts(t) \, dt + d\multiCountingProc(t).
\end{equation}

\begin{prop}\label{prop.compensatorClosedForm}
 Let $\multiCountingProc$ be an exponential-kernel Hawkes process started at
 $\multiCountingProc(0) = 0$.
 Then for every $\timeHorizon > 0$ and every type $e$,
 \begin{equation}\label{eq.compensatorClosedForm}
  \compensator_e(\timeHorizon)
  = \baseIntensity_e \timeHorizon
  + \sum_{\eone=1}^{\numEventTypes} \frac{\excitation\subscriptee}{\decay}
  \Big( \countingProc[\eone](\timeHorizon-) - \decayedCounts_{\eone}(\timeHorizon) \Big).
 \end{equation}
\end{prop}
\begin{proof}
 By \eqref{eq.intensityFromState} it suffices to compute
 $\int_{0}^{\timeHorizon}\decayedCounts_{\eone}$.
 The summands of \eqref{eq.decayedCounts} are non-negative,
 so Tonelli's theorem licenses the interchange, and
 \begin{equation*}
  \int_{0}^{\timeHorizon} \decayedCounts_{\eone}(u) \, du
  = \sum_{j \, : \, \arrivalTimes[\eone]_{j} < \timeHorizon}
    \int_{\arrivalTimes[\eone]_{j}}^{\timeHorizon} e^{-\decay(u - \arrivalTimes[\eone]_{j})} du
  = \sum_{j \, : \, \arrivalTimes[\eone]_{j} < \timeHorizon}
    \frac{1 - e^{-\decay(\timeHorizon - \arrivalTimes[\eone]_{j})}}{\decay} .
 \end{equation*}
 The sum runs over the arrivals of type $\eone$ \emph{strictly} before $\timeHorizon$,
 by the strict inequality of \eqref{eq.decayedCounts};
 there are $\countingProc[\eone](\timeHorizon-)$ of them,
 and their exponentials sum to $\decayedCounts_{\eone}(\timeHorizon)$.
\end{proof}

The compensator is thus available exactly, at any time, from the state and the counts,
with no numerical integration.

\begin{remark}\label{remark.readPreJump}
 Notice that the count in \eqref{eq.compensatorClosedForm} is read \emph{before}
 $\timeHorizon$ and the state \emph{at} it.
 The map
 $\timeHorizon \mapsto \countingProc[\eone](\timeHorizon-) -
 \decayedCounts_{\eone}(\timeHorizon)$
 is continuous, as it must be for Theorem \ref{thm.meyer1971} to apply:
 at an arrival of type $\eone$ both terms gain one.
 The post-jump reading
 $\countingProc[\eone](\timeHorizon) - \decayedCounts_{\eone}(\timeHorizon)$
 is not the compensator:
 it overstates \eqref{eq.compensatorClosedForm} by $\excitation\subscriptee/\decay$ at each
 arrival of type $\eone$, and agrees with it everywhere else.
 With $\decay = 1$, a single arrival at $t = 1/2$ and $\timeHorizon = 1/2$,
 the integral is $0$ while the post-jump reading gives $1$;
 and arrival times are precisely where Section \ref{sec.simulation} evaluates the compensator.
 It is also why we compute the compensator independently of the simulator rather than
 accumulating it along the path.
\end{remark}

The state moves deterministically between two arrivals, by $\varphi_t(z) = e^{-\decay t}z$,
and jumps at each of them.

\begin{thm}\label{thm.markovState}
 Let $\excitation \geq 0$, $\decay > 0$ and $\baseIntensity \geq 0$.
 Then $\decayedCounts$ is a piecewise-deterministic Markov process on
 $\R_{+}^{\numEventTypes}$,
 with flow $\varphi_t(z) = e^{-\decay t}z$,
 jump rate $\totalIntensity(z) = \mathbf{1}^{\transpose}(\baseIntensity + \excitation z)$
 and jump kernel
 $\sum_{e} \frac{\intensity(z)}{\totalIntensity(z)} \delta_{z + \unitVector_e}$,
 arbitrary where the rate vanishes.
  The extended generator of $\decayedCounts$ is
 \begin{equation}\label{eq.generatorOfDecayedCount}
  \generatorL f(z)
  = -\decay \, z \cdot \gradient f(z)
  + \sum_{e=1}^{\numEventTypes} \intensity(z)\big( f(z + \unitVector_e) - f(z) \big),
  \qquad \intensity(z) = \baseIntensity_e + (\excitation z)_e ,
 \end{equation}
 on every $f \in C^{1}(\R^{\numEventTypes})$ such that
 $\Expectation \int_{0}^{t} \sum_{e} \intensity(\decayedCounts(s))
 \lvert f(\decayedCounts(s) + \unitVector_e) - f(\decayedCounts(s)) \rvert \, ds < \infty$
 for every $t > 0$.
 Furthermore,
 the pair $(\multiCountingProc, \decayedCounts)$ is Markov,
 with generator
 \begin{equation}\label{eq.generatorOfHawkesProcess}
 \generatorL f(n,z) = -\decay \, z \cdot \gradient_{z} f(n,z)
 + \sum_{e} \intensity(z)\big( f(n + \unitVector_e, z + \unitVector_e) - f(n,z) \big) .
 \end{equation}
\end{thm}
\begin{proof}
 The intensity \eqref{eq.intensityFromState}, the jump rate and the jump kernel are functions
 of $\decayedCounts$ alone, so $\decayedCounts$ is Markov and $\multiCountingProc$ counts its
 jumps by type.
 Between arrivals $f(\decayedCounts)$ moves by the chain rule along the flow,
 at the rate $-\decay \, z \cdot \gradient f(z)$,
 which is the first term of \eqref{eq.generatorOfDecayedCount}.
 At an arrival of type $e$ it gains
 $f(\decayedCounts + \unitVector_e) - f(\decayedCounts)$,
 and by Definition \ref{def.compensator} such arrivals occur at the rate
 $\intensity(\decayedCounts)$,
 so the accumulated jumps minus the integral of the second term of
 \eqref{eq.generatorOfDecayedCount} is a local martingale.
 Adding the two gives \eqref{eq.generatorOfDecayedCount} up to localisation,
 which the stated integrability removes.
 The pair differs only in that $n$ does not flow,
 which is \eqref{eq.generatorOfHawkesProcess}.
\end{proof}
For the general statement on piecewise-deterministic Markov processes and for the domain of
the extended generator, see \cite[Theorem 5.5]{Dav84pie};
the Markov property of the exponential-kernel case is due to \cite{Oak75mar}.

\begin{remark}\label{remark.noSubcriticality}
 Notice that Theorem \ref{thm.markovState} assumes nothing about $\branchingRatio$.
 The integrability it requires holds on every finite horizon at every $\branchingRatio$,
 both for $f(z) = z$ and for $f(z) = z z^{\transpose}$.
 What supplies it is that for the exponential kernel the $k$-th convolution power
 $\hawkesKernel^{*k}$ carries $t^{k-1}/(k-1)!$,
 so $\sum_{k \geq 1}\int_{0}^{\timeHorizon}\hawkesKernel^{*k}$ converges at every
 $\branchingRatio$, and with it $\Expectation[\multiCountingProc(\timeHorizon)]$ and
 $\Expectation[\lvert\decayedCounts(\timeHorizon)\rvert^{2}]$.
 What subcriticality buys is stated in Section \ref{sec.stability}.
\end{remark}

\begin{exercise}\label{exercise.eventDependentDecays}
 A single $\decay$ shared by all pairs buys a state of dimension $\numEventTypes$.
 A decay $\decay\subscriptee$ per pair would require one decayed count per pair,
 a state of dimension $\numEventTypes^2$,
 and every expression of Sections \ref{sec.stability} and \ref{sec.secondOrder} would have to
 be rewritten on it.
 Rewrite them.
\end{exercise}

From now on, unless otherwise stated,
every Hawkes process is an exponential-kernel Hawkes process.

